%% Master CV - Pham Hai Nam
%% Comprehensive overview of all competencies, experience, and projects.
%% Use as a master reference when tailoring role-specific CVs (cv/main_<company>.tex).
%%
%% Populated by /setup on 2026-07-15.
%% Compile with: lualatex main_example.tex
%% (pdflatex often fails on modern MiKTeX with fontawesome5 font-expansion errors.)

\documentclass[11pt,a4paper,sans]{moderncv}
\moderncvstyle{banking}
\moderncvcolor{blue}

% Force the name and section headings to render in moderncv blue (color1).
% Default banking leaves them black: moderncvstylebanking.sty's \colorlet
% copies (not aliases) the pre-scheme accent colour, so the name colours are
% frozen before \moderncvcolor runs. Re-let them after. \namefont is the hook
% every name-style macro routes through, so this also works on moderncv 2.3.1
% (Debian/Ubuntu apt), which has no \firstnamestyle/\lastnamestyle at all.
\renewcommand*{\namefont}{\fontsize{34}{36}\bfseries\upshape}
\colorlet{firstnamecolor}{color1}
\colorlet{lastnamecolor}{color1}
\colorlet{namecolor}{color1}
\renewcommand*{\sectionstyle}[1]{{\sectionfont\color{color1}#1}}

\usepackage[utf8]{inputenc}
\usepackage{needspace}
% moderncv loads hyperref itself in an \AtEndPreamble hook, so \hypersetup
% must go in an \AtEndPreamble of our own: on moderncv < 2.4 a top-level
% \usepackage{hyperref} clashes with the class's own
% \RequirePackage[unicode]{hyperref}. From 2.4.0 the class passes its options
% through \PassOptionsToPackage instead, which is what removes that clash.
\AtEndPreamble{\hypersetup{
    colorlinks=true,
    linkcolor=blue,
    filecolor=magenta,
    urlcolor=blue,
    pdftitle={Pham Hai Nam - CV},
    % Keep pdfpagemode=UseNone: this block runs after moderncv's own
    % \AtEndPreamble (moderncv.cls sets pdfpagemode there), so a FullScreen
    % value here would win and open every CV in fullscreen presentation mode.
    pdfpagemode=UseNone,
}}
\usepackage[scale=0.80]{geometry}
\usepackage{import}

% personal data
\name{Pham Hai}{Nam}
\address{Ha Noi, Vietnam}{}{}
\phone[mobile]{0346294259}
\email{namphamhai7@gmail.com}
\extrainfo{\href{https://github.com/phamnam2003}{github.com/phamnam2003} {} \textbar{} {} \href{https://www.linkedin.com/in/pham-nam-153ab9259/}{linkedin.com/in/pham-nam-153ab9259}}

\begin{document}

\makecvtitle

% ============================================================
%     PROFILE STATEMENT
% ============================================================

\vspace{6pt}
\small{Fullstack developer with around three years across frontend and backend, strongest on Go backend services --- Gin, gRPC, Worker Pool, and dependency injection with uber-go/dig and Fx. Proposes and owns the database modeling and backend stack decisions for the modules he is responsible for, and builds the event-driven pipelines behind them with Kafka, RabbitMQ, and Redis, along with the observability and CI/CD around them. Also ships the ReactJS, Next.js, and Vue.js frontends that consume those APIs. Experience spans fintech and banking, government and public sector, real-estate CRM, and geospatial platforms.}

% ============================================================
%     CORE COMPETENCIES
% ============================================================

\section{Core Competencies}
\vspace{1pt}
\begin{itemize}

\item \textbf{Backend}: Go (Gin, gRPC with mutual TLS, interceptors and streaming, Worker Pool, dependency injection with uber-go/dig and Fx) and Node.js (Express, Strapi); RESTful APIs, WebSocket, Server-Sent Events; background workers and cron jobs; microservices.

\item \textbf{Frontend}: TypeScript/JavaScript with ReactJS, Next.js, Vue.js and component-based architecture; Redux, Redux Toolkit, Recoil, TanStack Query; Tailwind CSS, Ant Design, Shadcn UI.

\item \textbf{Distributed \& event-driven}: Apache Kafka and RabbitMQ messaging, Redis (cache, pub/sub, streams), Asynq, Ristretto; asynchronous import/export pipelines and background processing.

\item \textbf{Data}: PostgreSQL (incl. PostGIS), MySQL, Oracle, SQLite, MongoDB, ScyllaDB/Cassandra; database modeling, transaction management, sqlc; S3-compatible object storage (MinIO, RustFS, SeaweedFS).

\item \textbf{Infrastructure \& observability}: Docker, Kubernetes (Calico, Cilium, Envoy Gateway, Nginx Ingress); OpenTelemetry, Prometheus, Grafana, Loki; GitHub Actions, self-hosted GitLab Runners, ArgoCD, Helm; Linux (Ubuntu, Arch, CentOS).

\item \textbf{Testing, security \& AI-assisted development}: unit and integration tests across REST APIs, WebSocket, cron jobs and database interactions in Go and Node.js; JWT/PASETO, TOTP/2FA, TLS/mTLS; Python (Scrapy, Selenium) for e-commerce data crawling; Claude Code, GitHub Copilot, Spec-Kit.

\end{itemize}

% ============================================================
%     PROFESSIONAL EXPERIENCE
% ============================================================

\section{Professional Experience}
\vspace{3pt}
\begin{itemize}

% --- Most Recent Role ---
\needspace{5\baselineskip}
\item{\cventry{11/2025 - Present}{Frontend and Backend Developer}{AIONtech}{Ha Noi, Vietnam}{}{\vspace{1pt}
\begin{itemize}
    \item Implemented dependency injection with uber-go/dig and Fx to improve scalability and maintainability of Go backend services.
    \item Refactored and optimized internal packages to improve performance and strengthen application safety.
    \item Built background processing with Redis Pub/Sub and Apache Kafka for event-driven business workflows.
    \item Pushed realtime updates to clients over WebSocket and Server-Sent Events; built background task workers and scheduled cron jobs for recurring and long-running operations.
\end{itemize}}}

\vspace{3pt}

% --- Previous Role ---
\needspace{5\baselineskip}
\item{\cventry{06/2024 - 10/2025}{Frontend and Backend Developer}{Leeon Group}{Ha Noi, Vietnam}{}{\vspace{1pt}
\begin{itemize}
    \item Took ownership of an inherited production Go codebase from departing team members; maintained, debugged, and enhanced projects wired into CI/CD.
    \item Implemented messaging and caching: RabbitMQ, Redis (cache/pub-sub/streams), Ristretto, and the Worker Pool pattern for background tasks; worked with gRPC in depth (mutual TLS, metadata, interceptors, streaming).
    \item Stood up the observability stack: Prometheus, Grafana, Loki and OpenTelemetry --- traces, metrics, logs.
    \item Built ReactJS web applications with Tailwind CSS and Ant Design.
\end{itemize}}}

\vspace{3pt}

% --- Earlier Role ---
\needspace{5\baselineskip}
\item{\cventry{07/2023 - 02/2024}{Full-stack Developer}{TLGEO}{Ha Noi, Vietnam}{}{\vspace{1pt}
\begin{itemize}
    \item Built UIs with Vue.js and Next.js; developed and consumed APIs with ExpressJS and Strapi.
    \item Contributed to system design using PostgreSQL with the PostGIS extension for spatial data; used Mapbox for government mapping and agriculture projects; deployed on Linux/Ubuntu Server with SSH and Nginx.
\end{itemize}}}

\vspace{3pt}

% --- Oldest Role ---
\needspace{5\baselineskip}
\item{\cventry{04/2023 - 06/2023}{Frontend Intern}{Lalasoft}{Ha Noi, Vietnam}{}{\vspace{1pt}
\begin{itemize}
    % One printed line, deliberately. This is the lowest-signal bullet on the CV - a
    % three-month 2023 internship - and the words trimmed out of it ("ongoing",
    % "small", "tools as") were the cheapest line available on page 2. Page 2 has a
    % hard height ceiling beyond the 2-page rule: see the \enlargethispage note below.
    \item Applied Redux Core, Redux Toolkit, and Ant Design to company projects; built internal Chrome Extensions.
\end{itemize}}}

\end{itemize}

% ============================================================
%     SELECTED PROJECTS
% ============================================================

\section{Selected Projects}
\vspace{3pt}
\begin{itemize}

\needspace{4\baselineskip}
% Project entries put the project name in the third argument on purpose: moderncv's
% banking style renders argument 3 in bold and argument 2 in italics, so the project
% name has to sit third to be the line a reader scans first.
\item{\cventry{03/2026 - Present}{Backend Developer, team of 7}{SkyReality --- Real-estate lead platform}{AIONtech}{}{\vspace{1pt}
\begin{itemize}
    \item Proposed the lead and campaign data model plus the backend stack for the modules I owned; the design was reviewed and approved by the team before build.
    \item Integrated webhook event processing from Zalo, Slack, and Telegram bots; implemented OAuth2 Google Sign-In and multipart upload with checksum verification.
    % The realtime WebSocket/SSE + workers + cron bullet that sat here was cut as
    % redundant: the same capability is already claimed in the STM and C06 entries
    % below and in the AIONtech role bullets above. Its two lines are what buys the
    % Open-Source Projects section its place inside the 2-page budget.
    \item Built the platform's Next.js frontend (Shadcn UI, TanStack Query); deployed to Kubernetes with ArgoCD.
    \item Stack: Go (Gin), Kafka, Redis, PostgreSQL, uber-go/dig, MinIO; Kubernetes, ArgoCD; Next.js.
\end{itemize}}}

\vspace{3pt}

\needspace{4\baselineskip}
\item{\cventry{11/2025 - 01/2026}{Backend Developer, team of 8}{Sacombank Smart Teller Machine (STM)}{AIONtech}{}{\vspace{1pt}
\begin{itemize}
    \item Diagnosed a failing export (OOM kills and timeouts) as unbounded query result sets plus duplicate rows; paginated the query path and de-duplicated the output, cutting the same dataset from 167.8 MB to 16 MB.
    \item Built async import/export pipelines, step-level transaction logging, realtime WebSocket/SSE updates and cron-scheduled jobs; applied dependency injection for testability, and built the admin frontend in Next.js.
    \item Stack: Go (Gin), Kafka, Oracle, SQLite, Prometheus, Grafana, Docker; Next.js.
\end{itemize}}}

\vspace{3pt}

\needspace{4\baselineskip}
\item{\cventry{12/2025 - Present}{Backend Developer, team of 9}{Document AI for C06, Ministry of Public Security}{AIONtech}{}{\vspace{1pt}
\begin{itemize}
    \item Built and structured the core backend modules of this document-summarization and task-generation platform, contributing the data model and stack choices for that scope; now being generalized for customers beyond C06.
    \item Implemented realtime client updates over WebSocket and SSE, with background workers and cron jobs for long-running processing; deployed to Kubernetes with ArgoCD.
    \item Stack: Go (Gin), Kafka, PostgreSQL, Redis, MinIO, Docker, Kubernetes, ArgoCD.
\end{itemize}}}

\end{itemize}

% ============================================================
%     OPEN-SOURCE PROJECTS
% ============================================================

% Repos are named but not individually hyperlinked, matching the Core Competencies
% bullet this section replaced: the header already links the GitHub profile, and a
% per-repo link is a 404 risk the profile link does not carry. Entries are kept to
% one printed line each - this section is what pushes the CV against its 2-page
% budget, so a second wrapped line here costs a real line elsewhere.
\section{Open-Source Projects}
\vspace{1pt}
\begin{itemize}
\item \textbf{go-http-server/grpc} --- gRPC from the protocol up: four method types, interceptors, mTLS.
\item \textbf{phamnam2003/challenges} --- 20 Gang-of-Four patterns in Go, plus Kafka and ScyllaDB deep dives.
\item \textbf{go-http-server/temp} --- Gin template: PASETO, Asynq queue, sqlc, transactional writes, CI.
\end{itemize}

% ============================================================
%     EDUCATION
% ============================================================

% Stretch page 2 by a few lines so Education through References do not spill onto a
% page 3 of their own. The standard near-miss rescue per 05-cv-templates.md -
% compressing the geometry or the \vspace values instead is explicitly ruled out
% there. Placement matters: \enlargethispage only affects the page current when it
% is issued, so this has to sit ahead of the Education heading. Left further down
% (at Awards, where it used to live) it enlarges the very page 3 it should prevent.
%
% Page 2 has a second, quieter ceiling on top of the 2-page rule, and it bites before
% page 3 does. moderncv writes the \label{lastpage} behind the "n/m" footer from inside
% the footer box on the final page. Let page-2 content descend into that box - which is
% exactly what this \enlargethispage invites it to do - and the label is never written;
% \pageref{lastpage} then comes back undefined and the footer disappears from *every*
% page, page 1 included. It fails silently: no warning, no error, and the numbers are
% gone from the text layer as well, so the PDF still reads as fine.
%
% The trigger is how far page-2 content actually reaches, not this number on its own,
% so raising or lowering it is not a reliable fix - 4 with a full page 2 fails just as
% 5 does. Buy the room with a content cut, then check by grepping the extraction for
% both footers, not by eyeballing the PDF:
%   pdftotext -enc UTF-8 main_example.pdf - | grep -oE "\b[12]/2\b"   -> 1/2 and 2/2
\enlargethispage{4\baselineskip}
\section{Education \& Certifications}
\vspace{1pt}
\begin{itemize}

\needspace{3\baselineskip}
\item{\cventry{2021 - 2026}{Engineer's degree, Information Technology}{Hanoi Open University (HOU)}{Ha Noi, Vietnam}{}{\vspace{1pt}
Graduated 2026. Full-time programme, studied while working professionally from early 2023, with a concentrated study block in spring 2024. Thesis: an event-ticketing platform built for the sales-open traffic spike --- virtual waiting room ahead of the booking path, TTL-based distributed seat locking, Redis inventory cache.
}}
% The generic coursework list that used to close this entry (software engineering,
% data structures, databases, networks, web development) was cut to buy the thesis
% line its room. Relevance-weighted: for a developer with three years of shipped
% backend work, named coursework is the lowest-signal content on the page, while the
% thesis is a system-design story an interviewer can actually open a conversation on.
% Claim the design only - the platform was never load-tested, so no throughput,
% concurrency, or user-count figure may ever be added to this line.

% Certifications share the Education heading rather than taking one of their own. A
% moderncv section heading costs more vertical space than the line beneath it, and this
% page is already running on an \enlargethispage rescue - the same trade already made
% for "Languages, Interests & References" below. The heading still carries the word an
% ATS scans for. No issue dates exist for any of these, so none are stated.
%
% Short forms, not the registered titles. In full they are "S210: Using ScyllaDB
% Drivers", "S301: ScyllaDB Operations", "Alternator Course Completion Certificate",
% "ScyllaDB Labs 2025 Completion Certificate" and "LIVE Fall 2024 Advanced Track
% Completion Certificate" - see 01-candidate-profile.md, which holds the verbatim list
% for application forms. Written out here they run to two lines, and the second line is
% the one that puts this CV on a third page. The S210/S301 codes are the official
% identifiers, so a verifier can still find all five. "2025" and "Fall 2024" are parts
% of event names, not issue dates - never rewrite either into a date column.
\item \textbf{ScyllaDB} --- S210 Drivers; S301 Operations; Alternator; Labs 2025; LIVE Fall 2024 Advanced Track.

\end{itemize}

% ============================================================
%     AWARDS
% ============================================================

\section{Awards}
\vspace{1pt}
\begin{itemize}
% "Hanoi Open University" is deliberately not repeated here - it is on the scholarship
% line directly below and on the Education entry above, and dropping the third copy is
% what pulls this award back onto a single line.
\item \textbf{2nd prize, Faculty-level Student Scientific Research} (2024) --- topic: Agile in software development.
\item \textbf{Encouragement (merit) scholarship}, Hanoi Open University, awarded for three terms.
\end{itemize}

% ============================================================
%     LANGUAGES & REFERENCES
% ============================================================

% Languages, Interests and References share one heading on purpose. Each was a
% section holding a single line, and a moderncv section heading costs far more
% vertical space than the line under it - merging them reclaimed the room the
% Open-Source Projects section needed without dropping a word of any of them.
% Split them back into separate sections if a tailored copy has the space.
%
% Rust belongs on this line and nowhere else. It is active self-study with nothing
% shipped, so it must never appear in Core Competencies, where a reader would
% reasonably interview against it. The "nothing shipped" clause is not optional -
% see 01-candidate-profile.md, "Learning". It rides on the References line rather
% than taking an \item of its own purely for space; give it its own line first if a
% tailored copy has the room.
\section{Languages, Interests \& References}
\vspace{1pt}
\begin{itemize}
\item Vietnamese (native); English (professional working --- technical reading and writing).
\item Currently learning Rust (self-study, nothing shipped). References available upon request.
\end{itemize}

\end{document}
